\documentclass[10pt,a4paper]{article}
\usepackage{fontspec,amsmath,amssymb,graphicx}
\usepackage[margin=18mm]{geometry}
\setmainfont{Times New Roman}
\setlength{\parindent}{0pt}
\setlength{\parskip}{5pt}
\allowdisplaybreaks
\sloppy
\emergencystretch=2em
\begin{document}
\section*{BT\_3\_5\_3\_07 — Lời giải bổ sung}
\subsection*{Kiến thức và trình tự giải}
\textbf{Bước 1 — Chọn năng lượng theo toán tử.} Với sóng dùng $E=\frac12\int(u_t^2+a^2u_x^2)dx$. Với Klein–Gordon có thêm $u$, thêm $\frac12\int u^2dx$. Với dầm $u_{tt}=-a^2u_{xxxx}$ dùng $E=\frac12\int(u_t^2+a^2u_{xx}^2)dx$, không thay $u_{xx}$ bằng $u_x$.

\textbf{Bước 2 — Đạo hàm và tích phân từng phần.} Với $u_{tt}-a^2u_{xx}+\gamma u_t+\kappa u=F$, năng lượng có hạng $\kappa u^2$ cho


\[
E'=a^2[u_tu_x]_0^L-\gamma\int u_t^2dx+\int Fu_tdx.
\]


Với dầm, tích phân hai lần cho $E'=a^2[u_{tx}u_{xx}-u_tu_{xxx}]_0^L$. Viết riêng từng hạng biên rồi dùng điều kiện biên và đạo hàm theo thời gian của điều kiện đó để cho chúng bằng không.

\textbf{Bước 3 — Suy hệ quả.} Không cản, không nguồn thì $E$ bảo toàn. Có cản $\gamma>0$ thì $E'\le0$; nếu $E$ không có hệ số $1/2$, $E'=-2\gamma\int u_t^2\ge-2\gamma E$, suy $e^{-2\gamma t}E(0)\le E(t)\le E(0)$. Với hiệu hai nghiệm cùng dữ liệu, $E(0)=0$ nên có duy nhất.

\textbf{Bước 4 — Tách biến nếu phải tìm nghiệm cụ thể.} Sóng có cản $\gamma=1$ cho $T''+T'+a^2\lambda T=0$. Đặt $T=e^{-t/2}Y$, thu $Y''+(a^2\lambda-1/4)Y=0$, $Y(0)=A$, $Y'(0)=B+A/2$. Phân biệt dao động, nghiệm kép và hàm hyperbolic. Dầm cho $X''''=k^4X$, $T''+a^2k^4T=0$; dùng đủ bốn điều kiện biên để tìm phương trình của $kL$.

\textbf{Bước 5 — Kiểm tra hệ số và khẳng định của đề.} Năng lượng đầu chứa cả $g^2$ và các đạo hàm của $f$. Thử $f=0$, $g\ne0$ để phát hiện công thức bỏ vận tốc. Với dầm, nghiệm $k=0$ phải giải trực tiếp bằng đa thức bậc ba; nghiệm $0$ của phương trình lượng giác–hyperbolic có thể chỉ là nghiệm giả do cơ sở suy biến.

\textbf{Chi tiết phép triệt tiêu tích phân.} Đạo hàm động năng cho $\int u_tu_{tt}$; dùng PDE thay $u_{tt}$. Với sóng, $\int u_tu_{xx}=[u_tu_x]-\int u_{tx}u_x$, và tích phân cuối triệt tiêu đạo hàm thế năng $\int u_xu_{xt}$. Với dầm,


\[
\int u_tu_{xxxx}=[u_tu_{xxx}-u_{tx}u_{xx}]+\int u_{txx}u_{xx}.
\]


Tích phân cuối triệt tiêu đạo hàm $\frac12\int u_{xx}^2$ sau khi nhân hệ số và giữ đúng dấu PDE. Viết hạng biên trước, sau đó mới dùng điều kiện của từng đầu.

\textbf{Bước 1 — Tách biến và chọn bài toán trị riêng theo biên.} Đặt $u=X(x)T(t)$ trong $u_{tt}=a^2u_{xx}$, suy $X''+\lambda X=0$, $T''+a^2\lambda T=0$. Tích phân từng phần cho $\lambda\int X^2=\int (X')^2$ khi hạng biên bằng không; do đó $\lambda\ge0$. Giải cả trường hợp $\lambda=0$ trước khi xét $\lambda=k^2>0$.

\textbf{Bước 2 — Tìm hệ hàm riêng.} Trên $(0,L)$: Dirichlet hai đầu dùng $\sin(n\pi x/L)$, $n\ge1$; Neumann hai đầu dùng $1$ và $\cos(n\pi x/L)$; Dirichlet–Neumann dùng $\sin((n+1/2)\pi x/L)$, $n\ge0$; Neumann–Dirichlet dùng cos của các tần số nửa nguyên ấy. Chuẩn bình phương của mode khác không là $L/2$, của mode hằng là $L$.

\textbf{Bước 3 — Chiếu cả hai dữ liệu.} Với mode khác không đặt


\[
A_n=\frac2L\int_0^L fX_n\,dx,\quad B_n=\frac2L\int_0^L gX_n\,dx,\quad \omega_n=ak_n.
\]


Tích phân từng khoảng nếu dữ liệu chia đoạn; dùng tích phân từng phần hoặc công thức tích–tổng khi là đa thức/lượng giác. Hệ số $B_n$ là hệ số của vận tốc, nên hệ số đứng trước $\sin\omega_nt$ phải là $B_n/\omega_n$.

\textbf{Bước 4 — Ghép nghiệm và xét mode hằng.}


\[
u=\sum_n\left(A_n\cos\omega_nt+\frac{B_n}{\omega_n}\sin\omega_nt\right)X_n.
\]


Với Neumann hai đầu, thêm $A_0+B_0t$, trong đó $A_0=L^{-1}\int f$, $B_0=L^{-1}\int g$. Với tần số bị thay đổi bởi hạng $u$ hoặc cản $u_t$, giải lại ODE của từng mode; không dùng nguyên tần số $ak_n$.

\textbf{Bước 5 — Kiểm tra biên, dữ liệu và hội tụ.} Mỗi mode thỏa biên; tại $t=0$ hai chuỗi cho khai triển của $f,g$. Nếu dữ liệu không tương thích ở góc, giải theo nghĩa năng lượng hoặc từng miền và nêu đường mất tính trơn; không đạo hàm một chuỗi yếu tùy ý. Chu kỳ được suy từ tất cả mode có hệ số khác không, kể cả mode hằng; hai tần số có tỷ số vô tỷ ngăn một chu kỳ chung.

\textbf{Chi tiết tính tích phân hệ số.} Với $k=n\pi/L$, hai nguyên hàm thường dùng là


\[
\int x\sin kx\,dx=-x\cos kx/k+\sin kx/k^2,\qquad
\int x\cos kx\,dx=x\sin kx/k+\cos kx/k^2.
\]


Với đa thức $P$, ghi cả hai lần tích phân:


\[
\int_0^L P\sin kx\,dx
=\left[-P\cos kx/k\right]_0^L+
\left[P'\sin kx/k^2\right]_0^L-
\frac1{k^2}\int_0^L P''\sin kx\,dx.
\]


Chỉ bỏ hạng biên sau khi thế $0,L$ và điều kiện thực sự của $P$. Nếu dữ liệu chia đoạn, áp dụng công thức trên từng đoạn, rồi cộng các hạng tại điểm chia. Cuối cùng nhân nghịch đảo chuẩn của hàm riêng; nếu hệ số là của vận tốc, chia thêm tần số khi ghép nghiệm. Với mẫu số bằng $0$ trong tích–tổng, quay lại tích phân gốc để tính riêng.
\subsection*{Lời giải áp dụng cho bài}
\textbf{Bước 1. Đạo hàm năng lượng.} Với $u_{tt}+\xi u_t=u_{xx}$, $\xi\ge0$, biên cố định, $E=\int_0^L(u_t^2+u_x^2)dx$. Nhân PDE với $2u_t$ và tích phân từng phần:
\[E'=2[u_tu_x]_0^L-2\xi\int_0^Lu_t^2dx=-2\xi\int_0^Lu_t^2dx.\]
Hạng biên hết vì đạo hàm thời gian của vết Dirichlet bằng $0$.

\textbf{Bước 2. Ước lượng trên.} Tích phân bình phương không âm, nên $E'\le0$ và $E(t)\le E(0)$.

\textbf{Bước 3. Ước lượng dưới.} $\int u_t^2\le E$, nên $E'\ge-2\xi E$. Do đó $(e^{2\xi t}E)'=e^{2\xi t}(E'+2\xi E)\ge0$. So với giá trị đầu:
\[e^{-2\xi t}E(0)\le E(t)\le E(0).\]
\textbf{Bước 4. Kiểm tra hệ số.} Vì định nghĩa $E$ không có $1/2$, hệ số đạo hàm là $2\xi$. Với $\xi=0$, hai cận bằng nhau và thu bảo toàn năng lượng. Với $E(0)=0$ suy nghiệm duy nhất bằng $0$ khi li độ đầu cũng bằng $0$. Các công thức đúng trên đoạn $(0,L)$ sau khi thống nhất ký hiệu cận phải của đề.\end{document}